%% abtex2-modelo-relatorio-tecnico.tex, v<VERSION> laurocesar
%% Copyright 2012-2015 by abnTeX2 group at http://www.abntex.net.br/ 
%%
%% This work may be distributed and/or modified under the
%% conditions of the LaTeX Project Public License, either version 1.3
%% of this license or (at your option) any later version.
%% The latest version of this license is in
%%   http://www.latex-project.org/lppl.txt
%% and version 1.3 or later is part of all distributions of LaTeX
%% version 2005/12/01 or later.
%%
%% This work has the LPPL maintenance status `maintained'.
%% 
%% The Current Maintainer of this work is the abnTeX2 team, led
%% by Lauro César Araujo. Further information are available on 
%% http://www.abntex.net.br/
%%
%% This work consists of the files abntex2-modelo-relatorio-tecnico.tex,
%% abntex2-modelo-include-comandos and abntex2-modelo-references.bib
%%

% ------------------------------------------------------------------------
% ------------------------------------------------------------------------
% abnTeX2: Modelo de Relatório Técnico/Acadêmico em conformidade com 
% ABNT NBR 10719:2011 Informação e documentação - Relatório técnico e/ou
% científico - Apresentação
% Adaptado por Daniel Saad Nogueira Nunes para uso no IFB Taguatinga.
% ------------------------------------------------------------------------ 
% ------------------------------------------------------------------------


% Como opção escolha 'bacharelado' ou 'licenciatura'
\documentclass[licenciatura]{pre-projeto-computacao}
% ---
% Informações de dados para CAPA e FOLHA DE ROSTO
% ---
\titulo{Fatoração LZ78 utilizando R-Index}
\autor{Arthur Andrade D'Olival}
% para feminino use \orientadora{Nome da orientadora}
%\orientadora{Maria Bonita}
\orientador{Daniel Saad Nogueira Nunes}

% para masculino use \coorientador{Nome do coorientador}
%\coorientador{Padre Cícero}
\coorientadora{ - }

% Linha de pesquisa conforme tabela de áreas do CNPq
\areadepesquisa{Compressão de Dados}
\local{Brasília, DF}
\data{2026	}

\begin{document}

\selectlanguage{brazil}
\frenchspacing 
\imprimircapa
\imprimirfolhaderosto



\section*{Contextualização}

De acordo com projeções da Statista, empresa alemã especializada em dados de mercado e referência mundial em indicadores da economia digital, o volume global de dados criados, capturados, copiados e consumidos deverá atingir aproximadamente 390 zettabytes até 2028. 
Esse crescimento é impulsionado principalmente pela expansão de serviços digitais e pela crescente geração de dados em áreas como ciência, indústria e entretenimento.
Como consequência, surgem desafios cada vez maiores relacionados ao armazenamento, processamento e transmissão eficiente dessas informações, exigindo soluções capazes de lidar com volumes massivos de dados.

Nesse cenário, a compressão de dados torna-se um pilar essencial para o gerenciamento eficiente da informação, sendo mais do que uma simples ferramenta para economizar espaço em disco, passando a desempenhar um papel central no processamento da informação, permitindo que grandes volumes de dados sejam manipulados de forma eficiente e transmitidos através de redes.
Um exemplo claro dessa importância é a questão de versionamento em grandes repositórios de código fonte, onde é inviável manter um sistema que mantivesse cópias inteiras de diretórios a cada atualização sem o uso de algoritmos de compressão rodando em segundo plano.
Da mesma maneira, no campo da bioinformática, o sequenciamento de nova geração, Next Generation Sequencing, gera milhões de leituras de fragmentos de DNA que precisam ser alinhadas a genomas de referência contendo bilhões de pares de bases.
Lidar com esse grande volume de dados genômicos e manipulá-los em memória RAM só é computacionalmente viável graças às técnicas de compressão sem perdas.

Podemos classificar as diferentes abordagens de compressão de dados em dois grupos principais. 
O primeiro é a compressão com perda, lossy, amplamente utilizada em mídias de áudio, imagens e vídeos. Nesses casos, o algoritmo descarta detalhes imperceptíveis ou pouco relevantes aos seres humanos, permitindo alcançar taxas de compressão elevadas.
O segundo refere-se à compressão sem perdas, lossless, que garante a integridade absoluta da informação, algo fundamental em domínios críticos, como em bancos de dados genômicos, onde perder um único bit de informação pode comprometer análises científicas e resultados computacionais.

Dentro do universo da compressão sem perdas, a família de algoritmos Lempel-Ziv, especialmente o LZ78, desempenha um papel fundamental ao comprimir dados através da construção dinâmica de um dicionário, substituindo ocorrências repetidas por referências. À medida que o dicionário cresce e se torna mais rico, a capacidade do algoritmo de substituir trechos longos por referências compactas aumenta, tornando o processo cada vez mais eficiente.
Paralelamente, a Transformada de Burrows-Wheeler, BWT, introduziu uma abordagem inovadora e reversível que reorganiza os caracteres de um texto, agrupando símbolos semelhantes em blocos contínuos, chamados de runs. A evolução e combinação dessas técnicas permitiu o desenvolvimento de estruturas de dados auto-indexadas compactas, como o FM-index.
Tais estruturas são revolucionárias porque comprimem o texto e, simultaneamente, permitem buscas rápidas por padrões sem a necessidade de descompressão prévia. 

Contudo, apesar da eficiência dos índices tradicionais, o processamento de coleções de dados massivamente repetitivas revelou limitações no escalonamento do FM-index.
Nesse cenário, o r-index surge como uma evolução crítica, oferecendo uma estrutura cujo tamanho é proporcional ao número de runs da BWT, e não ao tamanho total do texto, garantindo alta performance em bases de dados de alta redundância.
Diante dessa perspectiva, o presente trabalho propõe a implementação da fatoração LZ78 fundamentada na estrutura do r-index.

\section*{Proposta}

A proposta central deste trabalho é propor a implementação do algoritmo de fatoração LZ78 apoiado nas capacidades de busca eficiente do r-index. O trabalho visa explorar a sinergia entre a compressão baseada na construção de dicionários dinâmicos do LZ78 e as estruturas modernas de indexação comprimida de textos.

\section*{Justificativa}

A realização deste projeto justifica-se pela demanda crítica e crescente por eficiência computacional em sistemas que lidam com redundância em escala extrema. A compressão deixou de ser apenas um utilitário e se tornou pilar essencial para o processamento e gerenciamento eficiente da informação.

A necessidade de otimizar os custos de memória na fatoração e na busca de dados atinge diretamente áreas de enorme impacto social e científico. Na bioinformática, o sequenciamento de nova geração, Next Generation Sequencing, gera bilhões de leituras de fragmentos de DNA. Alinhar e gerenciar esses volumes colossais de dados genômicos em memória RAM só se tornou possível graças aos índices baseados em BWT, como o FM-index, presente em ferramentas como o Bowtie. 
Da mesma forma, na engenharia de software, grandes repositórios baseiam-se fortemente na compressão de arquivos altamente repetitivos para efetuar o controle de versão de forma eficiente.
	
Ao estudar e aplicar a estrutura do r-index sobre um algoritmo clássico como o LZ78, este trabalho contribui para o avanço no manuseio de dados altamente repetitivos, viabilizando processos mais sustentáveis em termos de hardware computacional e acelerando a velocidade de análise de informações críticas.

\section*{Proposta Metodológica Preliminar}

Para o desenvolvimento do presente trabalho, a metodologia adotada consistirá, inicialmente, em uma pesquisa bibliográfica de caráter exploratório.
Esta etapa inicial tem como objetivo aprofundar o referencial teórico a respeito da compressão de dados sem perdas, algoritmos universais de modelagem de dicionários e o funcionamento de estruturas de dados auto-indexadas.
O embasamento teórico fundamentar-se-á em publicações científicas e literaturas consolidadas da ciência da computação, apoiando-se fortemente nos conceitos delineados por autores seminais da área, tais como Lempel, Ziv e Welch na compressão baseada em dicionário, Burrows e Wheeler no rearranjo reversível de strings, bem como Ferragina e Manzini no desenvolvimento do FM-index e estruturas oportunistas.

Em um segundo momento, a pesquisa assumirá uma abordagem de natureza aplicada e quantitativa.
O procedimento técnico central envolverá o projeto lógico e a implementação computacional do algoritmo de fatoração LZ78 operando de forma nativa sobre o \textit{r-index}.
A construção dessa solução fará uso de conceitos avançados de estruturas de dados sucintas e dicionários indexáveis, como o emprego das operações de busca de \textit{rank} e \textit{select} sobre blocos de caracteres.
Esse desenho metodológico garantirá que os dados permaneçam comprimidos em um espaço próximo ao limite ótimo teórico da informação, permitindo a extração de padrões e a construção da fatoração sem a exigência de descompressão prévia do texto.

Por fim, para a validação do modelo computacional desenvolvido e a consolidação dos dados empíricos, serão conduzidos experimentos práticos controlados.
A abordagem algorítmica proposta será submetida a testes de estresse utilizando bases de dados padronizadas e massivamente repetitivas. 
A avaliação quantitativa consistirá em medir rigorosamente as variáveis de tempo de execução e o pico de consumo de memória RAM.

\bibliography{bibliografia}

\end{document}
